\documentclass[12pt]{article}
\usepackage{amsmath,amssymb,amsfonts}
\usepackage[margin=1in]{geometry}
\usepackage{bm}

\newcommand{\vect}[1]{\mathbf{#1}}
\newcommand{\ehat}[1]{\hat{\vect{e}}_{#1}}

\begin{document}

\section*{Problem 2}

\textbf{Problem:} A particle moves in the $x_1 x_2$-plane so that its position vector is given by
\[
\vect{r} = a \cos \omega t \; \ehat{1} + a \sin \omega t \; \ehat{2} \qquad (\omega = \text{const}).
\]
Show that:
\begin{itemize}
\item[(a)] the velocity $\vect{v}$ is perpendicular to $\vect{r}$;
\item[(b)] the acceleration is directed toward the origin; find its magnitude;
\item[(c)] $\vect{L} = \vect{r} \times m\vect{v}$ is a constant vector; find its magnitude.
\end{itemize}

\bigskip
\textbf{Solution:}

\subsection*{(a) The velocity is perpendicular to $\vect{r}$}

The velocity is
\[
\vect{v} = \frac{d\vect{r}}{dt} = -a\omega \sin\omega t \; \ehat{1} + a\omega \cos\omega t \; \ehat{2}.
\]

Compute the dot product:
\begin{align}
\vect{r} \cdot \vect{v} &= (a\cos\omega t)(-a\omega\sin\omega t) + (a\sin\omega t)(a\omega\cos\omega t) \\
&= -a^2\omega \sin\omega t \cos\omega t + a^2\omega \sin\omega t \cos\omega t \\
&= 0.
\end{align}

Since $\vect{r} \cdot \vect{v} = 0$, the velocity is perpendicular to the position vector. \hfill $\blacksquare$

\textit{Note:} This can also be seen from
\[
\frac{d}{dt}(\vect{r} \cdot \vect{r}) = 2\,\vect{r} \cdot \vect{v} = \frac{d}{dt}(a^2) = 0.
\]
Since $|\vect{r}|^2 = a^2$ is constant, $\vect{r} \cdot \vect{v} = 0$ automatically.

\subsection*{(b) The acceleration is directed toward the origin}

The acceleration is
\[
\vect{a}_{\text{acc}} = \frac{d\vect{v}}{dt} = -a\omega^2 \cos\omega t \; \ehat{1} - a\omega^2 \sin\omega t \; \ehat{2} = -\omega^2 \vect{r}.
\]

Since $\vect{a}_{\text{acc}} = -\omega^2 \vect{r}$, the acceleration points in the direction opposite to $\vect{r}$, i.e., \textbf{toward the origin}.

The magnitude is:
\[
|\vect{a}_{\text{acc}}| = \omega^2 |\vect{r}| = \boxed{a\omega^2}.
\]

\subsection*{(c) Angular momentum is constant}

\begin{align}
\vect{L} &= \vect{r} \times m\vect{v} \\
&= m \bigl( a\cos\omega t \; \ehat{1} + a\sin\omega t \; \ehat{2} \bigr) \times \bigl( -a\omega\sin\omega t \; \ehat{1} + a\omega\cos\omega t \; \ehat{2} \bigr).
\end{align}

Using $\ehat{1} \times \ehat{1} = 0$, $\ehat{2} \times \ehat{2} = 0$, $\ehat{1} \times \ehat{2} = \ehat{3}$, $\ehat{2} \times \ehat{1} = -\ehat{3}$:
\begin{align}
\vect{L} &= m \bigl[ a\cos\omega t \cdot a\omega\cos\omega t \;(\ehat{1}\times\ehat{2}) + a\sin\omega t \cdot (-a\omega\sin\omega t) \;(\ehat{2}\times\ehat{1}) \bigr] \\
&= m a^2\omega \bigl[ \cos^2\omega t + \sin^2\omega t \bigr] \ehat{3} \\
&= m a^2 \omega \; \ehat{3}.
\end{align}

This is independent of time, confirming that $\vect{L}$ is constant. Its magnitude is:
\[
\boxed{|\vect{L}| = m a^2 \omega}.
\]

The angular momentum points in the $\ehat{3}$ direction (perpendicular to the plane of motion), consistent with the particle moving in a circle of radius $a$ with angular velocity $\omega$.

\end{document}
